% =====================================================================
\chapter{Полётные режимы и глобальные переменные}
% =====================================================================
\chapterintro
  {что такое полётные режимы EdgeTX (не путать с режимами полётного контроллера!), как
   у них устроены триммеры и плавный переход, как использовать глобальные переменные и
   настраивать параметры прямо в полёте}
  {модель самолёта или планера (для практики)}

\section{Полётные режимы EdgeTX}

\textbf{Полётный режим} (\emph{flight mode}, FM) в EdgeTX --- это набор триммеров и
условий для строк входов и миксеров, который включается переключателем. Типичные примеры
для планера: «Старт» (закрылки слегка вниз, триммер тангажа), «Круиз», «Скорость»,
«Термик», «Посадка» (кроу-тормоз).

\begin{note}
Полётные режимы пульта \emph{не имеют отношения} к режимам Betaflight/INAV (Angle, Horizon,
Acro). Режимы контроллера выбираются просто переключателем на AUX-канале.
\end{note}

\section{Страница Flight Modes}

\mpath{MDL\sep Flight Modes} содержит режимы \sw{FM0}…\sw{FM8}:

\begin{center}\small
\renewcommand{\arraystretch}{1.2}
\begin{tabularx}{\linewidth}{@{}p{30mm}X@{}}
\toprule
\menu{Name} & имя режима, показывается на главном экране и может проговариваться голосом \\
\menu{Switch} & переключатель, включающий режим. У \sw{FM0} переключателя нет --- это режим
  по умолчанию \\
\menu{Trims} & для каждой оси: свой триммер у режима, или использовать триммер другого режима
  (\sw{FM0}), или триммер другого режима плюс собственная добавка \\
\menu{Fade in / Fade out} & время плавного перехода в режим и из него, в секундах \\
\bottomrule
\end{tabularx}
\end{center}

\subsection{Приоритет}

Если активны переключатели сразу нескольких режимов, выигрывает режим \emph{с меньшим номером}:
\sw{FM1} важнее \sw{FM2}, \sw{FM2} важнее \sw{FM3} и т.\,д. \sw{FM0} работает, только когда
не активен ни один другой. На этом строят приоритет «Посадка» над всеми остальными: ставят её
в \sw{FM1} на отдельный переключатель.

\subsection{Использование в входах и миксерах}

У каждой строки входа и миксера есть параметр \menu{Flight modes} --- галочки
\code{0 1 2 3 4 5 6 7 8}. Снятая галочка означает, что строка в этом режиме не работает.
Так в режиме «Посадка» можно включить строку «закрылки вниз» и строку «компенсация тангажа».

\section{Шестипозиционный переключатель Boxer}

Кнопки 1--6 под экраном Boxer образуют переключатель \sw{6POS} (в списке источников он может
называться \sw{6P}). Каждая кнопка --- отдельное положение: \sw{6P1}…\sw{6P6}. Это идеальный
выбор полётных режимов:

\begin{center}\small
\begin{tabular}{@{}llll@{}}
\toprule
Режим & Имя & Switch & Триммеры \\
\midrule
\sw{FM0} & Cruise & --- & собственные \\
\sw{FM1} & Launch & \sw{6P1} & собственные \\
\sw{FM2} & Speed & \sw{6P2} & от FM0 + добавка \\
\sw{FM3} & Thermal & \sw{6P3} & от FM0 + добавка \\
\sw{FM4} & Land & \sw{6P4} & от FM0 \\
\bottomrule
\end{tabular}
\end{center}

Если нужно просто передать 6 положений в полётный контроллер (например, 6 режимов ArduPilot),
используйте \sw{6POS} как источник миксера канала: каждая кнопка даст своё значение
от $-100$ до $+100$.

На TX12 шестипозиционного переключателя нет; режимы выбирают комбинацией двух тумблеров
(например, \sw{SB} и \sw{SC}) или через логические переключатели.

\section{Глобальные переменные}

\textbf{Глобальная переменная} (\sw{GV1}…\sw{GV9}) --- именованное число, которое можно
использовать вместо константы почти в любом поле: вес, смещение, дифференциал, экспонента,
параметр логического переключателя.

Главная особенность: \textbf{значение GV может быть разным в каждом полётном режиме}. Так
делают, например, разный дифференциал элеронов в режимах «Круиз» и «Термик».

\subsection{Страница Global Variables}

\mpath{MDL\sep Global Variables}: для каждой переменной --- имя, значения в режимах
\sw{FM0}…\sw{FM8} (или ссылка «как в FM0»), минимум, максимум, единица измерения,
показ всплывающего окна при изменении (\menu{Popup}).

\subsection{Как подставить GV}

Встаньте на числовое поле (например, \menu{Weight} в миксере), долгое \btn{ENTER} →
\menu{Use GVar} (или аналог), выберите \sw{GV1}. Вместо числа появится \code{GV1}.

\section{Настройка «на лету»}

\begin{recipe}{подбираем экспоненту крутилкой в полёте}
\begin{enumerate}
  \item Создайте \sw{GV1} с именем \code{Expo}, диапазон 0…60.
  \item Во входе \sw{I1} (Ail) задайте \menu{Curve} = \menu{Expo} и значение \code{GV1}.
  \item В \menu{Special Functions}: условие = \sw{SD↓} (или любой свободный
        переключатель), функция \menu{Adjust GV1}, источник \sw{S1}.
  \item Пока \sw{SD} внизу, крутилка \sw{S1} меняет экспоненту от минимума до максимума,
        пульт показывает значение. Верните \sw{SD} --- значение зафиксируется.
  \item После полёта посмотрите значение \sw{GV1} и при желании впишите его как константу.
\end{enumerate}
\end{recipe}

\begin{summary}
\begin{itemize}
  \item Полётный режим = свой набор триммеров + условие для строк входов и миксеров.
  \item При нескольких активных режимах побеждает режим с меньшим номером.
  \item На Boxer удобно выбирать режимы кнопками \sw{6POS}.
  \item \sw{GV} --- числа, зависящие от полётного режима; их можно менять крутилкой в полёте.
\end{itemize}
\end{summary}
